%%
%% fall-lecture13.tex
%% 
%% Made by Alex Nelson <pqnelson@gmail.com>
%% Login   <alex@lisp>
%% 
%% Started on  2026-10-27T09:14:41-0700
%% Last update 2026-10-27T09:14:41-0700
%% 

\lecture{}

We will finish discussing nondimensionalization, then  pivot and
discuss Thermodynamics (over the next couple weeks).

\begin{example}[Drag on a falling sphere]
A sphere of diameter $d$ is falling and reaches some terminal velocity
$U_{\infty}$. The sphere is in equilibrium\footnote{Otherwise, it
would be experiencing some net force, and then it would experience
some nonzero acceleration and it would not reach terminal velocity.},
so the forces acting on it are the drag force $D$ and Earth's
gravitational pull $g$.

From the sphere's reference frame, the sphere is not
moving and the atmosphere is moving past it at velocity
$U_{\infty}$. The reference length scale is $d$ (the diameter of the
sphere) and the reference velocity is the terminal velocity
$U_{\infty}$. The Reynolds number is
\begin{equation}
\Rey = \frac{\rho U_{\infty}d}{\mu}=\frac{U_{\infty}d}{\nu}.
\end{equation}
The nondimensional drag must look like (c.f, the Drag Coefficient in Eq~\eqref{eq:lec12:me101a:drag-coef}):
\begin{equation}
\frac{D}{\rho U_{\infty}^{2}d^{2}}=f(\Rey)
\end{equation}
which is a function of the Reynolds number.
\begin{quote}
\textsc{Aside} (on body force): we have gravity acting on the sphere
(since the sphere is falling). We have
$\vec{\gamma}=-\rho\gradient{V}$ which can be absorbed into the
pressure since
\begin{equation}
-\partial_{i}p+\gamma_{i}=-\partial_{i}[p+V]=-\partial_{i}p_{\text{eff}}
\end{equation}
where $p_{\text{eff}}:=p+V$ is the ``effective pressure''.
\end{quote}
We can say something more about $f(\Rey)$. For Stokes flow (where $\Rey\ll1$)
\begin{equation}
f(\Rey)=\frac{C}{\Rey}
\end{equation}
where $C\in\RR$ is some [nonzero] constant which depends on the shape
of the body experiencing drag. Then the drag looks like
\begin{equation}
D=C\mu U_{\infty}d,
\end{equation}
which you should have seen before. (For a sphere, $C=6\pi$.)

Now things are in terminal velocity, what can be said of forces? As
noted earlier, the sphere is in equilibrium, so drag must equate to
the gravitational forcec
\begin{equation}
6\pi\mu U_{\infty}d  = \overbrace{\underbrace{\frac{4}{3}\pi\left(\frac{d}{2}\right)^{3}(\rho_{\text{solid}}-\rho_{\text{fluid}})}_{\text{effective mass}}\cdot{}g}^{\text{buoyancy}},
\end{equation}
which implies the terminal velocity must be
\begin{subequations}
  \begin{align}
U_{\infty}
&= \frac{1}{18}\cdot\frac{1}{2}\frac{d^{2}}{\mu}(\rho_{\text{solid}}-\rho_{\text{fluid}})g\\
&= \frac{2}{9}\cdot\frac{1}{8}\frac{d^{2}}{\nu}\left(\frac{\rho_{\text{solid}}}{\rho_{\text{fluid}}}-1\right).
  \end{align}
\end{subequations}
\end{example}

\begin{example}[Sand falling in water]
In this case, $\rho_{\text{solid}}/\rho_{\text{fluid}}\approx2$ and $\nu_{\text{fluid}}\approx0.0100~\mathrm{cm}^{2}/\mathrm{s}$.
We assumed $\Rey\ll1$, in this case
\begin{equation}
\Rey=\frac{U_{\infty}d}{\nu}\ll1\implies d\ll0.006~\mathrm{cm}=6\times10^{-5}~\mathrm{m},
\end{equation}
so the diameter of each grain of sand must be extraordinarily
fine. For sand larger than this, we are not in the Stokes regime.
[According to the Unified Soil Classification System, sand is
  classified as having a diameter between $7.4\times10^{-5}~\mathrm{m}$
and $4.75\times10^{-3}~\mathrm{m}$]
\end{example}

\begin{example}[Water in air]
[A.k.a., ``raindrops'']
In this case,
\begin{equation}
\frac{\rho_{\text{solid}}}{\rho_{\text{fluid}}}\approx780,
\end{equation}
and
\begin{equation}
\nu_{\text{air}}\approx0.15~\mathrm{cm}^{2}/\mathrm{s},
\end{equation}
and so
\begin{equation}
\Rey\ll1\implies d\ll0.004~\mathrm{cm}=4\times10^{5}~\mathrm{m}.
\end{equation}
This is the regime where we can use Stokes flow approximation.
[Wikipedia claims raindrops have sizes ranging from $1\times10^{-4}~\mathrm{m}$
to $9\times10^{-3}~\mathrm{m}$ in mean diameter.]
\end{example}

\begin{node}
A lot of nondimensional numbers we introduced are related to force
balance (the Mach number being the most famous exception). Here we
give one example of a nondimensional number not related to the balance
of forces.

Consider the convection diffusion equation for a passive scalar
\begin{equation}
\partial_{t}^{*}C+u_{i}^{*}\partial_{i}^{*}C=D(\partial_{i}^{*})^{2}C
\end{equation}
where $C$ is the passive scalar, the left-hand side describes the
convection, and the right-hand side describes the diffusion. We also
stress the use of dimensionful quantities by including the asterisks.

The first step is to stipulate the characteristic length scale $L$ and
velocity scale $U$, which permits us to write
\begin{equation}
u_{i}=\frac{u_{i}^{*}}{U},\quad\mbox{and}\quad x_{i}=\frac{x_{i}^{*}}{L}.
\end{equation}
Note that $[D]=\mathsf{L}^{2}\mathsf{T}^{-1}$. Then
\begin{equation}
\frac{U}{L}(\partial_{t}C+u_{i}\partial_{i}C)=\frac{U}{L}\left(\frac{D}{UL}\partial_{i}^{2}C\right).
\end{equation}
Dividing both sides by $U/L$ (since it's a nonzero quantity) gives us
\begin{equation}
\partial_{t}C+u_{i}\partial_{i}C = \frac{D}{UL}\partial_{i}^{2}C.
\end{equation}
We introduce the \define{Peclet Number} $\Peclet$ defined as
\begin{subequations}
  \begin{align}
\Peclet & := \frac{UL}{D}\\
&= \frac{UL}{\nu}\frac{\nu}{D}\\
&= (\Rey)(\Schmidt)
  \end{align}
\end{subequations}
where
\begin{equation}
\Schmidt := \frac{\nu}{D} = \frac{\mbox{diffusion of momentum}}{\mbox{diffusion of scalar}}
\end{equation}
is the \define{Schmidt Number}. This gives us an interesting
relationship between three nondimensional numbers.


Note: for the heat diffusion in Air, the Schmidt number is $\Schmidt\sim0.7$.
For the heat diffusion in Water, $\Schmidt\approx7$.
\end{node}

\subsection{Entropy Equations}

\begin{node}
This is just a more useful form of the energy equations.
\end{node}

\begin{recall}
The first Law of Thermodynamics 
\begin{equation}
\D e = \D Q + \D W
\end{equation}
where $e$ is the internal energy, $Q$ is the heat, and $\D W$ is the
work done \emph{on the system} by the Environment. (This is different
than how it is normally presented in thermodynamics, where we include
the work done \emph{by the system} on the Environment, which is why
the sign is flipped here.)
\end{recall}

\begin{recall}
The second Law of Thermodynamics takes the form, for reversible processes,
\begin{equation}
\D s = \frac{\D Q}{T}
\end{equation}
where $s$ is the specific entropy.
\end{recall}

\begin{recall}
The Gibbs relation
\begin{equation}
T\,\D s = \D e + p\,\D(1/\rho).
\end{equation}
We can think of $p\,\D(1/\rho)$ as describing compression/expansion,
and $T\,\D s$ as describing thermal/entropy contributions.
\end{recall}

\begin{goal}
We want to convert our Energy equation to be in terms of entropy.
\end{goal}

\begin{node}
We recall the internal energy equation:
\begin{align}
\rho\frac{D e}{D t}
&=-p\partial_{i}u_{i} + \tau_{ij}\partial_{j}u_{i} - \partial_{i}q_{i}
+ \rho Q\label{eq:me101a:fall:lec13:eq1}\\
&=\begin{pmatrix}\mbox{Pressure}\\
\mbox{Work}
\end{pmatrix} + (\mbox{dissipation}) - \begin{pmatrix}\mbox{Heat}\\
\mbox{Transfer}
\end{pmatrix} + \begin{pmatrix}\mbox{Heat}\\
\mbox{Source}
\end{pmatrix}\nonumber
\end{align}
where $\Phi:=\tau_{ij}\partial_{j}u_{i}$ is defined as the \define{Dissipation}.
\end{node}

\begin{node}
Now, we manipulate the continuity equation. Observe, by the chain rule,
\begin{equation}
\frac{D(\rho^{-1})}{Dt}=\frac{-1}{\rho^{2}}\frac{D\rho}{Dt}.
\end{equation}
Then
\begin{subequations}
  \begin{align}
\rho p\frac{D(1/\rho)}{Dt}
&=\frac{-p}{\rho}\frac{D\rho}{Dt}\quad\mbox{by chain rule}\\
&=\rho p(-\rho^{-2})(-\rho\partial_{i}u_{i})\quad\mbox{by continuity equation}\\
&=p\partial_{i}u_{i}.
  \end{align}
\end{subequations}
Hence
\begin{equation}\label{eq:me101a:fall:lec13:eq2}
\boxed{\rho p\frac{D(1/\rho)}{Dt}
=p\partial_{i}u_{i}.}
\end{equation}
\end{node}

\begin{node}
Now, we add Eq~\eqref{eq:me101a:fall:lec13:eq2} to Eq~\eqref{eq:me101a:fall:lec13:eq1}
to give us
\begin{equation}
\rho\left(\frac{De}{Dt}+p\frac{D(1/\rho)}{Dt}\right)=\Phi-\partial_{i}q_{i}+\rho Q.
\end{equation}
We use Gibbs relation to see the left-hand side is just $\rho T\,(D s/D t)$.
Therefore we obtain the \define{Entropy Equation}:
\begin{equation}
\boxed{\rho T\frac{D s}{D t}=\Phi-\partial_{i}q_{i}+\rho Q.}
\end{equation}
\end{node}

\begin{assume}
There is no internal heat source $Q=0$.
\end{assume}

\begin{assume}
Newtonian constitutive relation:
\begin{subequations}
  \begin{align}
\Phi &= \tau_{ij}\partial_{j}u_{i}\\
&= 2\mu(S_{ij}-\frac{1}{3}S_{kk}\delta_{ij})^{2}+\mu_{V}S_{kk}^{2}.
  \end{align}
\end{subequations}
In particular, $\Phi\geq0$.
\end{assume}

\begin{assume}
Fourier heat Law $q_{i}=-\kappa\partial_{i}T$.
\end{assume}

\begin{node}
We can use the continuity equation to infer
\begin{equation}
\partial_{t}(\rho s) = \rho\partial_{t}s + s\partial_{t}\rho=\rho\partial_{t}s-s\partial_{j}(\rho u_{j}).
\end{equation}
Then these simplify the Entropy equation, after expanding out the material
derivative, to become:
\begin{subequations}
  \begin{align}
\partial_{t}(\rho s)
&=-s\partial_{j}(\rho u_{j})-\rho u_{j}\partial_{j}s + \frac{1}{T}(\Phi+\partial_{j}(\kappa\partial_{j}T))\\
&=-\partial_{j}(s\rho u_{j}) + \frac{1}{T}\Phi + \frac{1}{T}\partial_{j}(\kappa\partial_{j}T).\label{eq:fall-lec13:me101a:rhs-to-be-studied}
  \end{align}
\end{subequations}
Let us examine how each of these three terms (on the right-hand side)
behave if we were to integrate them over the entire fluid.
\end{node}

\begin{assume}
Either
\begin{enumerate}
\item the volume is unbounded and as $|\vec{x}|\to\infty$ we have
both $u_{i}\to0$ and $T\to0$ sufficiently fast; or
\item the volume is bounded and $u_{i}=0$ and $n_{i}\partial_{i}T=0$ at the boundary.
\end{enumerate}
\end{assume}

\begin{node}
Then we see
\begin{subequations}
  \begin{align}
-\int_{V}\partial_{j}(s\rho u_{j})\,\D V &= -\oint_{\boundary V}\rho su_{j}n_{j}\,\D A\\
&=0\quad\mbox{by assumption}.
  \end{align}
\end{subequations}
So the first term in the right-hand side of Eq~\eqref{eq:fall-lec13:me101a:rhs-to-be-studied}
vanishes.
\end{node}

\begin{node}
By the Newtonian assumption, $\frac{1}{T}\Phi>0$. So the second term
in the right-hand side of Eq~\eqref{eq:fall-lec13:me101a:rhs-to-be-studied}
is positive.
\end{node}

\begin{node}
Integrating the third term in the right-hand side of Eq~\eqref{eq:fall-lec13:me101a:rhs-to-be-studied}:
\begin{subequations}
  \begin{align}
\int_{V}\frac{1}{T}\partial_{j}(\kappa\partial_{j}T)\,\D V
&=\int_{V}\partial_{j}\left(\frac{\kappa}{T}\partial_{j}T\right)\D V
+\int_{V}\frac{\kappa}{T^{2}}\partial_{j}^{2}T\,\D V\\
&=-\oint_{\boundary V}\frac{\kappa}{T}(\partial_{j}T)n_{j}\,\D A+\int_{V}\frac{\kappa}{T^{2}}\partial_{j}^{2}T\,\D V\\
&=0+\int_{V}\frac{\kappa}{T^{2}}\partial_{j}^{2}T\,\D V\quad\mbox{by boundary conditions}.
  \end{align}
\end{subequations}
We see that by assumption, this is a non-negative quantity.
\end{node}

\begin{node}
Combining all these results together, we can rewrite Eq~\eqref{eq:fall-lec13:me101a:rhs-to-be-studied}
in integral form as:
\begin{equation}
\partial_{t}\int_{V}\rho s\,\D V =
\int_{V}\frac{\kappa}{T^{2}}\partial_{j}^{2}T\,\D V+\int_{V}\frac{\Phi}{T}\,\D V.
\end{equation}
The right-hand side is the sum of a strictly positive and a
non-negative quantity, which means the right-hand side is strictly
positive. Taking a step back, this means that the change in entropy
over time (the left hand side) is strictly increasing. So, at the very
least, we recover familiar results (which is a good sanity check).
\end{node}
